% TOP_PROBLEM: {"problemId": "problem.polya-s-conjecture-for-dirichlet-eigenvalues", "canonicalTitle": "Pólya's Conjecture for Dirichlet Eigenvalues", "releaseRank": 402, "primaryDomain": "analysis_pde", "primaryDomainLabel": "Analysis & PDE", "displayStatus": "open_with_solved_subcases", "releaseStatus": "open_with_solved_subcases"}
% !TeX program = lualatex
% ATTEMPT_STATUS: unresolved
\documentclass[11pt]{article}
\usepackage[margin=25mm]{geometry}
\usepackage{fontspec}
\setmainfont{Latin Modern Roman}
\usepackage{amsmath,amssymb,mathtools}
\usepackage{unicode-math}
\setmathfont{Latin Modern Math}
\usepackage{xurl,hyperref}
\hypersetup{colorlinks=true,urlcolor=blue}
\setcounter{secnumdepth}{0}
\setlength{\parindent}{0pt}
\setlength{\parskip}{5pt}
\setlength{\emergencystretch}{3em}
\begin{document}
\section*{\texorpdfstring{Pólya's Conjecture for Dirichlet Eigenvalues}{Pólya's Conjecture for Dirichlet Eigenvalues}}\label{402}
\subsection{Definitions and mathematical statement}
Let $\Omega\subset{\mathbb{R}}^d$ be nonempty, bounded and open, $d\ge2$. The Dirichlet Laplacian is the nonnegative self-adjoint operator associated with the quadratic form $u\mapsto\int_\Omega|\nabla u|^2$ on $H_0^1(\Omega)$. List its eigenvalues increasingly with multiplicity as $\lambda_1\le\lambda_2\le\cdots$, and put $\omega_d=|\{x:\|x\|_2<1\}|$. P\'olya's assertion is
\[
 \lambda_k(\Omega)\ge(2\pi)^2
 \left(\frac{k}{\omega_d|\Omega|}\right)^{2/d}
 \qquad(k\ge1).
\]
No smoothness, connectedness, or tiling assumption on $\Omega$ is included. The inequality is for every eigenvalue, not just the leading asymptotic as $k\to\infty$. Neumann eigenvalues require a different formulation and inequality direction.
\par\smallskip\textit{Formulation and concept references:} \href{https://link.springer.com/article/10.1007/s00222-023-01198-1}{[S1]}.

\subsection{Short English statement}
Does every Dirichlet eigenvalue of every bounded Euclidean domain lie above its corresponding Weyl-law value?
\subsection{Research attempt}
\paragraph{Process and selected strategies.}
This attempt was generated by GPT-6 Astra Ultra on 27 September 2026. We prove the sharp pointwise inequality for rectangular boxes and their disjoint unions, derive the classical weaker universal estimate from Fourier mass, and test a proposed extension by approximation with boxes. The approximation fails in a precise spectral direction, illustrated by a domain whose area is unchanged when internal Dirichlet cuts are added.

The source [S1] proves the Dirichlet conjecture for Euclidean balls in every dimension, in addition to other stated cases. That is a genuine solved subcase, but it does not cover every bounded open set in the question. The averaged eigenvalue inequality of Li and Yau [E1] is also distinct from the required pointwise assertion. Our calculations keep this distinction explicit and do not assume smoothness, connectedness, or a tiling hypothesis for the unrestricted target.

\paragraph{Counting-function formulation and multiplicities.}
Write
\[
 C_d=(2\pi)^2\omega_d^{-2/d},\qquad
 N_\Omega(\Lambda)=\#\{j:\lambda_j(\Omega)\le\Lambda\}.
\]
The desired inequality is equivalent to
\[
 N_\Omega(\Lambda)\le(2\pi)^{-d}\omega_d|\Omega|\Lambda^{d/2}
 \quad\hbox{for all }\Lambda\ge0. \tag{402.1}
\]
If (402.1) holds, use $N_\Omega(\lambda_k)\ge k$ and rearrange. Conversely, if the eigenvalue inequality holds and $N_\Omega(\Lambda)=m\ge1$, then $\lambda_m\le\Lambda$ and the same rearrangement proves (402.1). Using $m$, rather than an arbitrarily chosen index with eigenvalue $\Lambda$, accounts for multiplicities.

For arbitrary bounded open $\Omega$, the $H_0^1$ form gives the usual compact Dirichlet resolvent. Indeed extend its functions by zero to a containing cube; the bounded $H_0^1$ unit ball is precompact in $L^2$ by compactness on the cube. This justifies a discrete eigenvalue sequence without imposing a regular boundary in the subsequent argument.

\paragraph{A sharp lattice count for boxes.}
For $Q=\prod_{i=1}^d(0,L_i)$, separation of variables gives eigenvalues
\[
 \lambda_m(Q)=\pi^2\sum_{i=1}^d(m_i/L_i)^2,
 \qquad m_i\in\mathbb Z_{\ge1}.
 \tag{402.2}
\]
For every tuple counted by $N_Q(\Lambda)$, attach the cell
\[
 B_m=\prod_i[(m_i-1)/L_i,m_i/L_i].
\]
Its volume is $1/\prod_iL_i$. The cells have disjoint interiors. Every coordinate of a point in $B_m$ is nonnegative and no greater than the corresponding coordinate of its upper corner, so the cell lies in the positive orthant of the ball of radius $\sqrt\Lambda/\pi$. Summing their volumes gives
\[
 \frac{N_Q(\Lambda)}{\prod_iL_i}
 \le2^{-d}\omega_d(\sqrt\Lambda/\pi)^d.
 \tag{402.3}
\]
Since $|Q|=\prod_iL_i$, this is exactly (402.1). The proof includes eigenvalues on the sphere: sharing cell faces does not affect volumes. It also works for unequal side lengths and in every dimension of the question.

\paragraph{Closure under disjoint unions.}
Suppose $\Omega$ is the disjoint union of open components $\Omega_j$, each satisfying (402.1), and their total volume is finite. The Dirichlet form is the orthogonal direct sum of the component forms, so the counting function is additive:
\[
 N_\Omega(\Lambda)=\sum_jN_{\Omega_j}(\Lambda)
 \le(2\pi)^{-d}\omega_d\Lambda^{d/2}\sum_j|\Omega_j|.
 \tag{402.4}
\]
For countably many components, the inequality follows first for finite partial sums and then by monotone convergence. The displayed finite upper bound also ensures that only finitely many eigenvalues can be counted at a fixed $\Lambda$. Thus any bounded disjoint union of boxes satisfies the desired assertion. The same reasoning combines any other established class of components, such as balls using [S1]. It does not glue adjacent components across their Dirichlet boundaries.

\paragraph{A universal Fourier estimate and the exact factor it loses.}
Let $\phi_1,\ldots,\phi_k$ be orthonormal Dirichlet eigenfunctions, extended by zero. Use the unitary Fourier transform and set
\[
 F(\xi)=\sum_{j=1}^k|\widehat\phi_j(\xi)|^2,
 \qquad M=(2\pi)^{-d}|\Omega|.
\]
Bessel's inequality applied to the exponential restricted to $\Omega$, followed by Plancherel, gives
\[
 0\le F\le M,\qquad \int F=k,\qquad
 \int |\xi|^2F(\xi)\,d\xi=\sum_{j=1}^k\lambda_j.
\]
Choose $R$ by $M\omega_dR^d=k$. Pointwise,
$(|\xi|^2-R^2)(F-M1_{\{|\xi|\le R\}})\ge0$.
Integration and equality of the two masses imply
\[
 \sum_{j=1}^k\lambda_j
 \ge M\int_{|\xi|\le R}|\xi|^2\,d\xi
 =\frac d{d+2}C_d|\Omega|^{-2/d}k^{1+2/d}.
 \tag{402.5}
\]
This is the Li--Yau lower bound [E1], reconstructed here with its constants. Since $\lambda_k$ is at least the average of the first $k$ eigenvalues, it gives
\[
 \lambda_k\ge\frac d{d+2}C_d(k/|\Omega|)^{2/d}.
 \tag{402.6}
\]
The factor $d/(d+2)<1$ remains. Summing eigenvalue bounds and bounding one individual eigenvalue are different operations; the average estimate cannot simply be relabeled as the pointwise conjecture.

The relaxed optimization in frequency space is exact for the constraints $0\le F\le M$ and $\int F=k$: filling the centered ball minimizes its second moment. Thus those constraints alone cannot produce a stronger sum inequality by the same minimization. Any improved use of this route must exploit additional information about Fourier transforms of eigenfunctions or control a counting function more directly. No such strengthening is established here.

\paragraph{Why approximation from inside has the wrong direction.}
One might partition most of an arbitrary domain into many small boxes, apply (402.3) to their union, and let their total volume approach $|\Omega|$. If $\Omega_{\rm in}\subseteq\Omega$, Dirichlet domain monotonicity gives
\[
 \lambda_k(\Omega_{\rm in})\ge\lambda_k(\Omega).
\]
A lower bound for the larger eigenvalue on the left does not imply a lower bound for the smaller one on the right. Equivalently, the inner domain has fewer low eigenvalues, so an upper bound for its counting function gives no upper bound for the counting function of $\Omega$.

An outer approximation would have the useful direction: if $\Omega\subseteq\Omega_j$ and $\Omega_j$ satisfies the conjecture, then
\[
 \lambda_k(\Omega)\ge C_d(k/|\Omega_j|)^{2/d}.
\]
If such outer domains could always be chosen with volumes tending to $|\Omega|$, the conclusion would follow. But disjoint open grid boxes generally omit their separating grid hyperplanes and hence fail to contain $\Omega$. Adding neighborhoods of those interfaces changes the domain into a connected or overlapping union for which the already-proved disjoint-union argument does not apply. The existence of suitable outer approximants is exactly the unproved step in this strategy.

\paragraph{A measure-zero change with a large spectral effect.}
In dimension two take the unit square $Q=(0,1)^2$ and remove the internal vertical segments at $x=j/m$, $1\le j<m$. The remaining open set $Q_m$ consists of $m$ strips of width $1/m$ and height one. It has the same area as $Q$, but (402.2) gives
\[
 \lambda_1(Q)=2\pi^2,\qquad
 \lambda_1(Q_m)=\pi^2(m^2+1). \tag{402.7}
\]
The first eigenvalue of the disconnected union is the minimum of the identical component first eigenvalues, so the second equality is exact. Thus equality, or convergence, of volumes is insufficient to justify convergence of Dirichlet spectra. The removed interfaces have area zero but enforce additional trace conditions in $H_0^1$. This explicitly defeats a box-approximation proof that treats the grid cuts as spectrally irrelevant merely because they are null sets.

Weyl asymptotics likewise give the sharp leading constant only as $k\to\infty$ for each fixed domain under the applicable formulation. The conjecture asks for a one-sided bound for every index. A relative error tending to zero does not determine its sign at all finite indices; an asymptotic statement cannot replace that extra inequality.

\paragraph{Verification and outcome.}
The finite checks count positive lattice points for selected boxes, verify multiplicities and the strip spectrum in (402.7), and independently simplify the constants in the volume argument. The Fourier and domain-monotonicity steps are justified analytically above. No numerical eigenvalue approximation is used to assert the unrestricted inequality.

\textbf{UNRESOLVED.} The sharp bound is proved for boxes and their disjoint unions, and the universal estimate (402.6) is recovered with its weaker factor. The outer-approximation route lacks suitable containing domains of asymptotically equal volume; no sharp pointwise inequality for arbitrary bounded open sets is established.

\subsection{Sources}
\begin{itemize}
\item[S1]N. Filonov, M. Levitin, I. Polterovich, and D. Sher, Inventiones Mathematicae 234 (2023), 129–169.
\url{https://link.springer.com/article/10.1007/s00222-023-01198-1}.
\textit{Formulation source.}
\item[E1]P. Li and S.-T. Yau, \emph{On the Schr\"odinger equation and the eigenvalue problem}, Communications in Mathematical Physics 88 (1983), 309--318.
\url{https://link.springer.com/article/10.1007/BF01213210}.
\textit{Averaged Dirichlet eigenvalue inequality; bibliographic record and theorem statement consulted 27 September 2026.}
\end{itemize}
\end{document}
